\chapter{What the Rule Costs}
\label{sec:discussion}

beni's design follows a principle: \emph{guarantees, not restrictions}. The language keeps Elm's
guarantees---no run-time exception, no silent wrong answer, exhaustive matches, managed effects---but a
rule that does not buy such a guarantee does not belong in the language as an error: it becomes a
warning, gets an escape hatch, or is dropped. A limit that bounds a real blow-up and can be lifted
by an annotation is the model of a restriction done right. Every rule in this design is to be tested
against that principle.

This rule buys a guarantee about \emph{cost}: \emph{an in-place edit the program asks for, by calling a
named writer, is performed in place}---\code{set}, \code{update}, \code{swap} and \code{pop} in $O(1)$,
\code{push} and \code{cons} in amortised $O(1)$, \code{append} in time proportional to what it appends,
\code{insertAt} and \code{removeAt} in the $O(n)$ element move a hand-written program also pays---and a
list that is built rather than edited costs what the same construction costs in JavaScript. There is
no hidden copy, no trie and no reference count, and the runtime is smaller. The price is that programs
that are \emph{correct} under value semantics are rejected. The design choice made here is that the
error is wanted. This chapter states its price without softening it.

\section{Programs that are rejected}
\label{sec:disc-rejected}

Every one of the following is a correct beni program today.
\begin{enumerate}
\item Keeping an old version and then editing it in place: an undo history, a snapshot for a diff, a
``previous'' field in the model, a test that holds its input and compares it with the output
(\S\ref{sec:hard-undo}).
\item A command, fiber, subscription or stored closure that captured a list which a later message edits
in place (\S\ref{sec:hard-fibers}).
\item Sending a list that the sender still holds: the beginner's
\code{Cmd.perform (λsend → send (Saved model.todos))} is rejected at the \code{send}, because sending
consumes the message while the model still holds \code{model.todos} (\S\ref{sec:hard-tea}, \Obl{P6}).
\item A closure that edits a captured list in place and is passed to \code{map}, \code{filter},
\code{sortBy} or any other $\mathit{many}$ position, or called twice (\S\ref{sec:hard-closures}).
\item An in-place edit of a list read out of a container that is still live: a dictionary lookup
followed by a push, or a push on an element of a list of lists whose contents are not exclusive
(\S\ref{sec:hard-containers}).
\item An in-place edit of a list from JavaScript, a \code{Ref}, a top-level value, \code{init} while it is
a value, or a \code{foreign} without a summary (\S\ref{sec:hard-foreign}, \S\ref{sec:hard-ref},
\S\ref{sec:hard-tea}).
\item An in-place edit of a rendered list on a platform that keeps what it rendered: on beni's current
template runtime, every list \code{view} shows (\S\ref{sec:hard-tea}).
\item A named writer whose two arguments are one list, such as \code{List.append xs xs}, which fails
\Disjoint; \code{xs ++ xs} builds a new list and is accepted.
\item Anything the checker's joins, path cut or caps cannot follow (\S\ref{sec:hard-branches}): limits.
\end{enumerate}
The fix is one call, \code{List.copy}, in every case, and the message writes it; where the edit need not
be in place, the non-writing form---\code{++}, a spread, \code{List.map}---is the other fix. As Henriksen
puts it for Futhark, ``an alias-related type error can always be fixed by copying''
\citep{henriksen2022uniqueness}. What the fix costs is a copy where the program meant to share. For the
first case that is the program's intended semantics made explicit; for the last it is a copy the
programmer pays for the checker's blindness.

\section{Programs that copy}
\label{sec:disc-copy}

List syntax builds new arrays (\S\ref{sec:model-demands}). Today beni's representation makes
\code{acc ++ [ x ]} and a prepend \code{[ x, …acc ]} amortised $O(1)$ by claiming a free slot at either
end of the trie; under the rule each copies \code{acc}. An accumulator written that way in a loop the
building-loop rewrite does not catch---a fold's lambda, a non-tail recursion---becomes quadratic, as it
would be in JavaScript and as \code{acc ++ [ x ]} already is in Elm. The checker knows when the named
writer would have been accepted, because it knows whether \code{acc} is unique and dead at the site,
and in exactly that case it warns: ``this copies \code{acc}; \code{List.push acc x} would write it in
place''. It is a warning, never an error, and never a change of meaning. The evaluation counts these
warnings, separately for loops and elsewhere.

\section{What it costs the representation}

The trie's reason to exist is gone, and so is its persistence: a program that wants versions pays
$O(n)$ per version, explicitly. The named prepend, \code{List.cons}, keeps amortised $O(1)$ on a unique
list through headroom in front of a view (\S\ref{sec:disc-options}, Change~8), at the cost of views
being more common than today and of the space that headroom takes. Without that change, a named
prepend on a unique plain array would be \code{unshift}, $O(n)$, and the guarantee above would have to
say so.

\section{What it costs the API}

A public function's summary is part of its contract. Editing a body can break a caller in another
module with an ownership error that the caller's author did not write---the usual objection to
inferring signatures, that changing the body of a function would change its signature. The interface
firewall makes this visible, not painless. It is also the objection Henriksen raises against
fresh-result annotations: they clutter interfaces, and the programmer ``will not be able to be
completely ignorant of this facet of the type system'' \citep{henriksen2026aliasing}. Inferred summaries
are not written, but they are published, and a caller depends on them; non-locality is, if anything,
greater than with written annotations, since the cause of an error may be a body the caller never
read.

\section{What it costs learners}

The error will fire first in \code{update}. Crichton and Krishnamurthi found the chapter on ownership
to be ``a significant drop-off point'' for readers of the Rust book \citep[\S3.1]{crichton2024profiling},
and in Crichton, Gray and Krishnamurthi's study learners who could usually say why a program was
rejected fixed it in only 46\,\% of cases \citep[\S2.4]{crichton2023grounded}. The message design of
Chapter~\ref{sec:errors} is the mitigation; it is not a guarantee, and the evaluation plan measures
something about people for this reason (\S\ref{sec:eval-people}).

\section{What two real applications say}
\label{sec:disc-apps}

We examined two applications written in beni: a TodoMVC implementation (in four builds) and an
implementation of Conduit, the RealWorld example application. Both run on beni's TEA platform over the
\emph{template} runtime, which keeps every list it rendered; neither runs on the direct platform. By our
reading of their sources, neither calls \code{set}, \code{push}, \code{update}, \code{swap}, \code{pop},
\code{insertAt} or \code{removeAt}; they edit lists with \code{List.map}, \code{List.filter} and
\code{++}. An earlier version of this analysis claimed that the rule as stated changed nothing in them
except to accept their concatenations; that claim was wrong in four ways, and we withdraw it:
\begin{itemize}
\item it analysed them on the direct platform, which they do not use, while their own platform keeps
every rendered list, so any in-place edit of \code{model.todos} or of a rendered error list would be an
error there;
\item TodoMVC's \code{init} is a top-level value, so every list that starts in it is escaped;
\item TodoMVC persists its list from a command, \code{save todos = Cmd.do λ⊤ → ⋯ (encode todos)}, which
encodes \emph{inside} the thunk, so the list is captured by a command that the runtime does not promise
to run before the next message;
\item its helper \code{changed model todos = \{ model | todos = todos \}} was argued to read nothing of
\code{.todos}, which that version's summary lattice could not express.
\end{itemize}

What is true under the revised rules is simpler. Because \code{++} builds a new list, neither
application contains a demand at all: TodoMVC's one concatenation per build, \code{model.todos ++ [ new ]},
and Conduit's sixteen (seven with a model field on the left, nine with a list built in the same
function) all copy, as they do in JavaScript. Both applications are \emph{accepted, on their own
platform, and nothing in them is written in place}. The rule costs them nothing and buys them nothing
beyond the smaller list runtime.

The in-place writes they could benefit from are their \code{List.map model.todos (λt → ⋯)} and
\code{List.filter} calls. Under Change~1, which adds in-place forms of these under their own names, a
TodoMVC rewritten to use them would be:
\begin{itemize}
\item \emph{rejected on its own platform}, because the template runtime keeps the rendered
\code{model.todos};
\item \emph{rejected on the direct platform as written}, because \code{init} is a value and
\code{save}'s command captures the list;
\item \emph{accepted on the direct platform} only with \code{init} a function (Change~12), with
\code{save} encoding before the command is made or with the platform draining sync commands
(\Obl{P7}), and with the helper's $\unused$ summary, which the revised lattice has
(Appendix~\ref{app:examples}).
\end{itemize}
The core library has seven calls of \code{List.push}, every one an accumulator in a fold or a loop: in
combining and partitioning results, in listing a dictionary's keys, values and pairs, and in a retry
schedule. The compiler's test programs for the direct platform have about fifty in-place edits of keyed
lists, written to exercise its edit scripts. So the code that would use the rule today is library code
and tests written for the direct platform; whether application code would use it is the evaluation's
question.

\section{The run-time alternative}
\label{sec:disc-runtime}

Henriksen's 2026 advice to other designers is to avoid this fight: ``unless you have a good reason to
pick this fight, don't do it'', because abstraction hides aliasing, fresh-result annotations clutter
interfaces and convenience functions produce ``spurious and annoying aliasing errors''; he suggests that
``most languages that want in-place updates should probably just use APL/SaC/Koka/Lean-style reference
counting, which reuses memory when the count is one and copies otherwise'', keeping a static discipline
for languages that ``must guarantee that no allocation takes place'' \citep{henriksen2026aliasing}. His
reasons are not specific to type systems, and our design being an inferred pass does not answer them:
he is himself exploring inference outside Futhark's language of types.

Our own prototype study supports his alternative more than it supports us. A run-time scheme with a
sticky ``shared'' bit on each array and no decrements made the same in-place decisions as a full
reference count in every scenario; our static variant matched its in-place writes only with the help
of inlined folds and per-field summaries (Appendix~\ref{app:method}). The bit accepts every program and never errs. Measured against it, the rule
buys three things and only three: \emph{a cost the programmer can read from the source}---an edit
written in place either is in place or is an error, never a silent $O(n)$ copy that turns a loop
quadratic, which in the same study cost a fold 8.2 times the baseline's time---\emph{a smaller runtime},
with no bit to test on every write and no copy-on-shared path to ship; and \emph{the absence of run-time
ownership state}, which in a garbage-collected target has no other use. Whether those are worth the
rejections is not a question the design can answer. It is why the evaluation compares the two
(\S\ref{sec:eval-decisions}) and why the warning mode below exists.

\section{Escape hatches}
\label{sec:disc-hatches}

\code{List.copy}, always. The non-writing forms---\code{++}, spreads, \code{List.map}---which never
demand anything. A persistent-vector library written in beni for programs that want versions, not a
representation of \code{List}. And a \emph{warning mode}---a flag such as \code{--ownership=warn}---that
keeps a run-time fallback (the sticky shared bit above) and copies where the checker would reject,
printing the same message as a warning. That is the principle's ``make it a warning'' route, and it is
the form every language with a run-time fallback ships in some way: Lean offers \code{dbgTraceIfShared},
which prints a message ``if the argument's reference count is greater than one'' \citep{lean2025rc}, and
Koka's \code{fip} check reports where a function is not fully in place \citep{lorenzen2023fp2}. But it
gives up the one thing the rule is for: an accepted program's cost would again depend on whether its
author read a warning. We do not propose it as the end state; we propose measuring first
(Chapter~\ref{sec:evaluation}) and deciding between error and warning on criteria fixed in advance.
